\documentclass[
  jou,
  floatsintext,
  longtable,
  a4paper,
  nolmodern,
  notxfonts,
  notimes,
  12pt,
  colorlinks=true,linkcolor=blue,citecolor=blue,urlcolor=blue]{apa7}

\usepackage{amsmath}
\usepackage{amssymb}



\usepackage[bidi=default]{babel}
\babelprovide[main,import]{spanish}


% get rid of language-specific shorthands (see #6817):
\let\LanguageShortHands\languageshorthands
\def\languageshorthands#1{}

\RequirePackage{longtable}
\RequirePackage{threeparttablex}

\makeatletter
\renewcommand{\paragraph}{\@startsection{paragraph}{4}{\parindent}%
	{0\baselineskip \@plus 0.2ex \@minus 0.2ex}%
	{-.5em}%
	{\normalfont\normalsize\bfseries\typesectitle}}

\renewcommand{\subparagraph}[1]{\@startsection{subparagraph}{5}{0.5em}%
	{0\baselineskip \@plus 0.2ex \@minus 0.2ex}%
	{-\z@\relax}%
	{\normalfont\normalsize\bfseries\itshape\hspace{\parindent}{#1}\textit{\addperi}}{\relax}}
\makeatother




\usepackage{longtable, booktabs, multirow, multicol, colortbl, hhline, caption, array, float, xpatch}
\usepackage{subcaption}


\renewcommand\thesubfigure{\Alph{subfigure}}
\setcounter{topnumber}{2}
\setcounter{bottomnumber}{2}
\setcounter{totalnumber}{4}
\renewcommand{\topfraction}{0.85}
\renewcommand{\bottomfraction}{0.85}
\renewcommand{\textfraction}{0.15}
\renewcommand{\floatpagefraction}{0.7}

\usepackage{tcolorbox}
\tcbuselibrary{listings,theorems, breakable, skins}
\usepackage{fontawesome5}

\definecolor{quarto-callout-color}{HTML}{909090}
\definecolor{quarto-callout-note-color}{HTML}{0758E5}
\definecolor{quarto-callout-important-color}{HTML}{CC1914}
\definecolor{quarto-callout-warning-color}{HTML}{EB9113}
\definecolor{quarto-callout-tip-color}{HTML}{00A047}
\definecolor{quarto-callout-caution-color}{HTML}{FC5300}
\definecolor{quarto-callout-color-frame}{HTML}{ACACAC}
\definecolor{quarto-callout-note-color-frame}{HTML}{4582EC}
\definecolor{quarto-callout-important-color-frame}{HTML}{D9534F}
\definecolor{quarto-callout-warning-color-frame}{HTML}{F0AD4E}
\definecolor{quarto-callout-tip-color-frame}{HTML}{02B875}
\definecolor{quarto-callout-caution-color-frame}{HTML}{FD7E14}

%\newlength\Oldarrayrulewidth
%\newlength\Oldtabcolsep


\usepackage{hyperref}



\usepackage{color}
\usepackage{fancyvrb}
\newcommand{\VerbBar}{|}
\newcommand{\VERB}{\Verb[commandchars=\\\{\}]}
\DefineVerbatimEnvironment{Highlighting}{Verbatim}{commandchars=\\\{\}}
% Add ',fontsize=\small' for more characters per line
\usepackage{framed}
\definecolor{shadecolor}{RGB}{241,243,245}
\newenvironment{Shaded}{\begin{snugshade}}{\end{snugshade}}
\newcommand{\AlertTok}[1]{\textcolor[rgb]{0.68,0.00,0.00}{#1}}
\newcommand{\AnnotationTok}[1]{\textcolor[rgb]{0.37,0.37,0.37}{#1}}
\newcommand{\AttributeTok}[1]{\textcolor[rgb]{0.40,0.45,0.13}{#1}}
\newcommand{\BaseNTok}[1]{\textcolor[rgb]{0.68,0.00,0.00}{#1}}
\newcommand{\BuiltInTok}[1]{\textcolor[rgb]{0.00,0.23,0.31}{#1}}
\newcommand{\CharTok}[1]{\textcolor[rgb]{0.13,0.47,0.30}{#1}}
\newcommand{\CommentTok}[1]{\textcolor[rgb]{0.37,0.37,0.37}{#1}}
\newcommand{\CommentVarTok}[1]{\textcolor[rgb]{0.37,0.37,0.37}{\textit{#1}}}
\newcommand{\ConstantTok}[1]{\textcolor[rgb]{0.56,0.35,0.01}{#1}}
\newcommand{\ControlFlowTok}[1]{\textcolor[rgb]{0.00,0.23,0.31}{\textbf{#1}}}
\newcommand{\DataTypeTok}[1]{\textcolor[rgb]{0.68,0.00,0.00}{#1}}
\newcommand{\DecValTok}[1]{\textcolor[rgb]{0.68,0.00,0.00}{#1}}
\newcommand{\DocumentationTok}[1]{\textcolor[rgb]{0.37,0.37,0.37}{\textit{#1}}}
\newcommand{\ErrorTok}[1]{\textcolor[rgb]{0.68,0.00,0.00}{#1}}
\newcommand{\ExtensionTok}[1]{\textcolor[rgb]{0.00,0.23,0.31}{#1}}
\newcommand{\FloatTok}[1]{\textcolor[rgb]{0.68,0.00,0.00}{#1}}
\newcommand{\FunctionTok}[1]{\textcolor[rgb]{0.28,0.35,0.67}{#1}}
\newcommand{\ImportTok}[1]{\textcolor[rgb]{0.00,0.46,0.62}{#1}}
\newcommand{\InformationTok}[1]{\textcolor[rgb]{0.37,0.37,0.37}{#1}}
\newcommand{\KeywordTok}[1]{\textcolor[rgb]{0.00,0.23,0.31}{\textbf{#1}}}
\newcommand{\NormalTok}[1]{\textcolor[rgb]{0.00,0.23,0.31}{#1}}
\newcommand{\OperatorTok}[1]{\textcolor[rgb]{0.37,0.37,0.37}{#1}}
\newcommand{\OtherTok}[1]{\textcolor[rgb]{0.00,0.23,0.31}{#1}}
\newcommand{\PreprocessorTok}[1]{\textcolor[rgb]{0.68,0.00,0.00}{#1}}
\newcommand{\RegionMarkerTok}[1]{\textcolor[rgb]{0.00,0.23,0.31}{#1}}
\newcommand{\SpecialCharTok}[1]{\textcolor[rgb]{0.37,0.37,0.37}{#1}}
\newcommand{\SpecialStringTok}[1]{\textcolor[rgb]{0.13,0.47,0.30}{#1}}
\newcommand{\StringTok}[1]{\textcolor[rgb]{0.13,0.47,0.30}{#1}}
\newcommand{\VariableTok}[1]{\textcolor[rgb]{0.07,0.07,0.07}{#1}}
\newcommand{\VerbatimStringTok}[1]{\textcolor[rgb]{0.13,0.47,0.30}{#1}}
\newcommand{\WarningTok}[1]{\textcolor[rgb]{0.37,0.37,0.37}{\textit{#1}}}

\providecommand{\tightlist}{%
  \setlength{\itemsep}{0pt}\setlength{\parskip}{0pt}}
\usepackage{longtable,booktabs,array}
\newcounter{none} % for unnumbered tables
\usepackage{calc} % for calculating minipage widths
% Correct order of tables after \paragraph or \subparagraph
\usepackage{etoolbox}
\makeatletter
\patchcmd\longtable{\par}{\if@noskipsec\mbox{}\fi\par}{}{}
\makeatother
% Allow footnotes in longtable head/foot
\IfFileExists{footnotehyper.sty}{\usepackage{footnotehyper}}{\usepackage{footnote}}
\makesavenoteenv{longtable}

\usepackage{graphicx}
\makeatletter
\newsavebox\pandoc@box
\newcommand*\pandocbounded[1]{% scales image to fit in text height/width
  \sbox\pandoc@box{#1}%
  \Gscale@div\@tempa{\textheight}{\dimexpr\ht\pandoc@box+\dp\pandoc@box\relax}%
  \Gscale@div\@tempb{\linewidth}{\wd\pandoc@box}%
  \ifdim\@tempb\p@<\@tempa\p@\let\@tempa\@tempb\fi% select the smaller of both
  \ifdim\@tempa\p@<\p@\scalebox{\@tempa}{\usebox\pandoc@box}%
  \else\usebox{\pandoc@box}%
  \fi%
}
% Set default figure placement to htbp
\def\fps@figure{htbp}
\makeatother







\usepackage{newtx}

\defaultfontfeatures{Scale=MatchLowercase}
\defaultfontfeatures[\rmfamily]{Ligatures=TeX,Scale=1}





\title{Introducción al entorno R}


\shorttitle{INTRO R PROGRAMACIÓN}


\usepackage{etoolbox}



\ccoppy{\textcopyright~2021}



\author{Edison Achalma}



\affiliation{
{Escuela Profesional de Economía, Universidad Nacional de San Cristóbal
de Huamanga}}




\leftheader{Achalma}

\date{2021-03-02}


\abstract{Este abstract será actualizado una vez que se complete el
contenido final del artículo. }

\keywords{keyword1, keyword2}

\authornote{\par{\addORCIDlink{Edison Achalma}{0000-0001-6996-3364}} 
\par{ }
\par{   El autor no tiene conflictos de interés que revelar.    Los
roles de autor se clasificaron utilizando la taxonomía de roles de
colaborador (CRediT; https://credit.niso.org/) de la siguiente
manera:  Edison Achalma:   conceptualización, metodología, análisis
formal, investigación, recursos, curación de
datos, redacción, visualización, supervisión, administración del
proyecto}
\par{La correspondencia relativa a este artículo debe dirigirse a Edison
Achalma, Escuela Profesional de Economía, Universidad Nacional de San
Cristóbal de Huamanga, Portal Independencia N°
57, Ayacucho, AYA 5001, Perú, Email: \href{mailto:elmer.achalma.09@unsch.edu.pe}{elmer.achalma.09@unsch.edu.pe}}
}

\usepackage{pbalance}
% \usepackage{float}
\makeatletter
\let\oldtpt\ThreePartTable
\let\endoldtpt\endThreePartTable
\def\ThreePartTable{\@ifnextchar[\ThreePartTable@i \ThreePartTable@ii}
\def\ThreePartTable@i[#1]{\begin{figure}[!htbp]
\onecolumn
\begin{minipage}{0.485\textwidth}
\oldtpt[#1]
}
\def\ThreePartTable@ii{\begin{figure}[!htbp]
\onecolumn
\begin{minipage}{0.48\textwidth}
\oldtpt
}
\def\endThreePartTable{
\endoldtpt
\end{minipage}
\twocolumn
\end{figure}}
\makeatother


\makeatletter
\let\endoldlt\endlongtable		
\def\endlongtable{
\hline
\endoldlt}
\makeatother

\newenvironment{twocolumntable}% environment name
{% begin code
\begin{table*}[!htbp]%
\onecolumn%
}%
{%
\twocolumn%
\end{table*}%
}% end code

\urlstyle{same}



\hypersetup{
  pdftitle={Introducci\'{o}n al entorno R},
  pdfauthor={Autor}
}
\makeatletter
\@ifpackageloaded{caption}{}{\usepackage{caption}}
\AtBeginDocument{%
\ifdefined\contentsname
  \renewcommand*\contentsname{Tabla de contenidos}
\else
  \newcommand\contentsname{Tabla de contenidos}
\fi
\ifdefined\listfigurename
  \renewcommand*\listfigurename{Lista de Figuras}
\else
  \newcommand\listfigurename{Lista de Figuras}
\fi
\ifdefined\listtablename
  \renewcommand*\listtablename{Lista de Tablas}
\else
  \newcommand\listtablename{Lista de Tablas}
\fi
\ifdefined\figurename
  \renewcommand*\figurename{Figura}
\else
  \newcommand\figurename{Figura}
\fi
\ifdefined\tablename
  \renewcommand*\tablename{Tabla}
\else
  \newcommand\tablename{Tabla}
\fi
}
\@ifpackageloaded{float}{}{\usepackage{float}}
\floatstyle{ruled}
\@ifundefined{c@chapter}{\newfloat{codelisting}{h}{lop}}{\newfloat{codelisting}{h}{lop}[chapter]}
\floatname{codelisting}{Listado}
\newcommand*\listoflistings{\listof{codelisting}{Lista de Listados}}
\makeatother
\makeatletter
\makeatother
\makeatletter
\@ifpackageloaded{caption}{}{\usepackage{caption}}
\@ifpackageloaded{subcaption}{}{\usepackage{subcaption}}
\makeatother
\makeatletter
\@ifpackageloaded{fontawesome5}{}{\usepackage{fontawesome5}}
\makeatother

% From https://tex.stackexchange.com/a/645996/211326
%%% apa7 doesn't want to add appendix section titles in the toc
%%% let's make it do it
\makeatletter
\xpatchcmd{\appendix}
  {\par}
  {\addcontentsline{toc}{section}{\@currentlabelname}\par}
  {}{}
\makeatother

%% Disable longtable counter
%% https://tex.stackexchange.com/a/248395/211326

\usepackage{etoolbox}

\makeatletter
\patchcmd{\LT@caption}
  {\bgroup}
  {\bgroup\global\LTpatch@captiontrue}
  {}{}
\patchcmd{\longtable}
  {\par}
  {\par\global\LTpatch@captionfalse}
  {}{}
\apptocmd{\endlongtable}
  {\ifLTpatch@caption\else\addtocounter{table}{-1}\fi}
  {}{}
\newif\ifLTpatch@caption
\makeatother

\begin{document}

\maketitle


\hypertarget{toc}{}
\tableofcontents
\newpage
\section[Introduction]{Introducción al entorno R}

\setcounter{secnumdepth}{3}

\setlength\LTleft{0pt}




\section{Guía completa de introducción a
R}\label{guuxeda-completa-de-introducciuxf3n-a-r}

\begin{quote}
Referencia técnica independiente sobre los \textbf{fundamentos del
lenguaje R}: qué es, cómo se usa de forma interactiva, su sintaxis
básica, objetos, modos, atributos, operadores, funciones y el sistema de
ayuda --- todo lo previo a la manipulación de datos propiamente dicha
(data frames, importación de datos, \texttt{dplyr}, etc., que se tratan
en una guía aparte). Basada en el manual oficial \textbf{``An
Introduction to R''} del R Core Team, publicado por CRAN
(cran.r-project.org/doc/manuals), no en blogs ni tutoriales de terceros.
\end{quote}

\begin{center}\rule{0.5\linewidth}{0.5pt}\end{center}

\subsection{Tabla de contenidos}\label{tabla-de-contenidos}

\begin{enumerate}
\def\labelenumi{\arabic{enumi}.}
\tightlist
\item
  \hyperref[1-quuxe9-es-r]{¿Qué es R?}
\item
  \hyperref[2-r-y-la-estaduxedstica-una-diferencia-de-filosofuxeda]{R y
  la estadística: una diferencia de filosofía}
\item
  \hyperref[3-usar-r-de-forma-interactiva]{Usar R de forma interactiva}
\item
  \hyperref[4-el-sistema-de-ayuda-de-r]{El sistema de ayuda de R}
\item
  \hyperref[5-sintaxis-buxe1sica-comandos-sensibilidad-a-mayuxfasculas-comentarios]{Sintaxis
  básica: comandos, sensibilidad a mayúsculas, comentarios}
\item
  \hyperref[6-asignaciuxf3n-de-variables]{Asignación de variables}
\item
  \hyperref[7-historial-de-comandos]{Historial de comandos}
\item
  \hyperref[8-ejecutar-comandos-desde-un-archivo-y-redirigir-salida]{Ejecutar
  comandos desde un archivo y redirigir salida}
\item
  \hyperref[9-permanencia-de-datos-y-el-espacio-de-trabajo]{Permanencia
  de datos y el espacio de trabajo}
\item
  \hyperref[10-vectores-la-estructura-fundamental]{Vectores: la
  estructura fundamental}
\item
  \hyperref[11-aritmuxe9tica-vectorizada-y-la-regla-de-reciclaje]{Aritmética
  vectorizada y la regla de reciclaje}
\item
  \hyperref[12-generaciuxf3n-de-secuencias]{Generación de secuencias}
\item
  \hyperref[13-vectores-luxf3gicos]{Vectores lógicos}
\item
  \hyperref[14-valores-faltantes-na-y-nan]{Valores faltantes:
  \texttt{NA} y \texttt{NaN}}
\item
  \hyperref[15-vectores-de-caracteres]{Vectores de caracteres}
\item
  \hyperref[16-vectores-de-uxedndices-seleccionar-y-modificar-subconjuntos]{Vectores
  de índices: seleccionar y modificar subconjuntos}
\item
  \hyperref[17-objetos-modo-longitud-y-atributos]{Objetos: modo,
  longitud y atributos}
\item
  \hyperref[18-coerciuxf3n-entre-modos]{Coerción entre modos}
\item
  \hyperref[19-la-clase-de-un-objeto]{La clase de un objeto}
\item
  \hyperref[20-estructuras-de-control-condicionales-y-bucles]{Estructuras
  de control: condicionales y bucles}
\item
  \hyperref[21-funciones-propias]{Funciones propias}
\item
  \hyperref[22-paquetes-el-sistema-de-extensiuxf3n-de-r]{Paquetes: el
  sistema de extensión de R}
\item
  \hyperref[23-espacios-de-nombres-namespaces]{Espacios de nombres
  (namespaces)}
\item
  \hyperref[24-buenas-pruxe1cticas-seguxfan-el-manual-oficial]{Buenas
  prácticas según el manual oficial}
\item
  \hyperref[25-referencias-oficiales]{Referencias oficiales}
\end{enumerate}

\begin{center}\rule{0.5\linewidth}{0.5pt}\end{center}

\subsection{1. ¿Qué es R?}\label{quuxe9-es-r}

Según la documentación oficial, R es un \textbf{conjunto integrado de
herramientas de software} para manipulación de datos, cálculo y
representación gráfica. Entre otras cosas, ofrece:

\begin{itemize}
\tightlist
\item
  una herramienta efectiva de manejo y almacenamiento de datos,
\item
  un conjunto de operadores para cálculos sobre arreglos, en particular
  matrices,
\item
  una colección amplia, coherente e integrada de herramientas
  intermedias para análisis de datos,
\item
  herramientas gráficas para análisis y visualización de datos, tanto en
  pantalla como en copia impresa,
\item
  un lenguaje de programación bien desarrollado, simple y efectivo
  (llamado \textbf{S}), que incluye condicionales, bucles, funciones
  recursivas definidas por el usuario, y funciones de entrada y salida.
\end{itemize}

El término ``entorno'' (\emph{environment}) pretende caracterizar a R
como un sistema plenamente planeado y coherente, en lugar de una
acumulación incremental de herramientas muy específicas e inflexibles,
como ocurre frecuentemente con otro software de análisis de datos.

R puede considerarse una implementación del lenguaje \textbf{S},
desarrollado originalmente en los Laboratorios Bell por Rick Becker,
John Chambers y Allan Wilks, y que también es la base de los sistemas
S-PLUS.

R es, sobre todo, un vehículo para el desarrollo de nuevos métodos de
análisis de datos interactivo. Se ha extendido rápidamente mediante una
gran colección de \textbf{paquetes} (\emph{packages}). La mayoría de los
programas escritos en R son, según el propio manual, esencialmente
efímeros: escritos para un único análisis de datos puntual.

\begin{center}\rule{0.5\linewidth}{0.5pt}\end{center}

\subsection{2. R y la estadística: una diferencia de
filosofía}\label{r-y-la-estaduxedstica-una-diferencia-de-filosofuxeda}

El manual oficial señala una diferencia importante de filosofía entre R
(y S) y otros sistemas estadísticos clásicos como SAS o SPSS. En estos
últimos, un análisis estadístico típicamente produce una salida
abundante de forma inmediata. En S (y por tanto en R), un análisis
estadístico normalmente se realiza como \textbf{una serie de pasos}, con
resultados intermedios almacenados en objetos: R produce una salida
mínima en cada paso y almacena los resultados en un objeto (por ejemplo,
un objeto de ajuste de un modelo) para su posterior interrogación
mediante otras funciones de R.

Esto tiene una implicación práctica importante para quien empieza: en R,
\emph{no ver} una salida extensa tras ajustar un modelo no significa que
algo haya fallado; significa que el resultado quedó guardado en un
objeto, listo para ser examinado con las funciones apropiadas
(\texttt{summary()}, \texttt{print()}, etc.).

R incluye unos 25 paquetes llamados ``estándar'' y ``recomendados'' que
vienen junto con cualquier instalación de R, y muchos más disponibles a
través de la red de sitios CRAN (cran.r-project.org).

\begin{center}\rule{0.5\linewidth}{0.5pt}\end{center}

\subsection{3. Usar R de forma
interactiva}\label{usar-r-de-forma-interactiva}

\subsubsection{3.1. El prompt}\label{el-prompt}

Cuando se usa el programa R, este emite un \emph{prompt} (símbolo de
espera de instrucción) cuando espera comandos de entrada. El prompt por
defecto es \texttt{\textgreater{}}. Si un comando no está completo al
final de una línea, R emite un prompt distinto por defecto:

\begin{verbatim}
+
\end{verbatim}

en la segunda línea y subsiguientes, y continúa leyendo hasta que el
comando esté sintácticamente completo.

\subsubsection{3.2. Procedimiento recomendado al usar R por primera vez
(Linux/Unix)}\label{procedimiento-recomendado-al-usar-r-por-primera-vez-linuxunix}

El manual recomienda, para una primera sesión bajo UNIX/Linux, el
siguiente procedimiento:

\textbf{Paso 1 --- Crear un subdirectorio de trabajo separado}, por
ejemplo llamado \texttt{work}, para contener los archivos de datos sobre
los que se trabajará en ese problema concreto:

\begin{Shaded}
\begin{Highlighting}[]
\ExtensionTok{$}\NormalTok{ mkdir work}
\ExtensionTok{$}\NormalTok{ cd work}
\end{Highlighting}
\end{Shaded}

\textbf{Paso 2 --- Iniciar el programa R}:

\begin{Shaded}
\begin{Highlighting}[]
\ExtensionTok{$}\NormalTok{ R}
\end{Highlighting}
\end{Shaded}

\textbf{Paso 3} --- A partir de aquí, pueden emitirse comandos de R.

\textbf{Paso 4 --- Salir del programa} con el comando:

\begin{Shaded}
\begin{Highlighting}[]
\SpecialCharTok{\textgreater{}} \FunctionTok{q}\NormalTok{()}
\end{Highlighting}
\end{Shaded}

En ese momento se preguntará si se desea guardar los datos de la sesión
(workspace). Según el sistema, esto aparecerá como un cuadro de diálogo
o como un prompt de texto al que se puede responder \texttt{yes},
\texttt{no} o \texttt{cancel} (basta una sola letra de abreviación) para
guardar los datos antes de salir, salir sin guardar, o volver a la
sesión. Los datos guardados estarán disponibles en futuras sesiones de
R.

Sesiones posteriores son simples: situarse en el directorio
\texttt{work} e iniciar el programa de nuevo:

\begin{Shaded}
\begin{Highlighting}[]
\ExtensionTok{$}\NormalTok{ cd work}
\ExtensionTok{$}\NormalTok{ R}
\end{Highlighting}
\end{Shaded}

\subsubsection{3.3. En Windows}\label{en-windows}

El procedimiento es básicamente el mismo: se crea una carpeta como
directorio de trabajo, y esta se establece en el campo ``Start In'' del
acceso directo de R. Luego se lanza R haciendo doble clic en el ícono.

\begin{quote}
\textbf{Nota práctica}: si usas RStudio (en cualquier sistema
operativo), el equivalente moderno y recomendado a ``crear una carpeta
de trabajo separada'' es usar \textbf{RStudio Projects} (\texttt{File} →
\texttt{New\ Project...}), que gestiona automáticamente el directorio de
trabajo, el historial y el workspace por proyecto.
\end{quote}

\begin{center}\rule{0.5\linewidth}{0.5pt}\end{center}

\subsection{4. El sistema de ayuda de R}\label{el-sistema-de-ayuda-de-r}

R incluye un sistema de ayuda integrado, similar a la utilidad
\texttt{man} de UNIX. Es, según el propio manual, la primera herramienta
a dominar antes de avanzar.

\subsubsection{4.1. Ayuda sobre una función
específica}\label{ayuda-sobre-una-funciuxf3n-especuxedfica}

Para obtener información sobre una función específica, por ejemplo
\texttt{solve}:

\begin{Shaded}
\begin{Highlighting}[]
\SpecialCharTok{\textgreater{}} \FunctionTok{help}\NormalTok{(solve)}
\end{Highlighting}
\end{Shaded}

Forma alternativa, más breve:

\begin{Shaded}
\begin{Highlighting}[]
\SpecialCharTok{\textgreater{}}\NormalTok{ ?solve}
\end{Highlighting}
\end{Shaded}

\subsubsection{4.2. Ayuda sobre operadores o palabras con significado
sintáctico}\label{ayuda-sobre-operadores-o-palabras-con-significado-sintuxe1ctico}

Para una característica especificada por caracteres especiales, el
argumento debe ir entre comillas dobles o simples, convirtiéndolo en una
``cadena de caracteres'' (\emph{character string}). Esto también es
necesario para algunas palabras con significado sintáctico, incluyendo
\texttt{if}, \texttt{for} y \texttt{function}:

\begin{Shaded}
\begin{Highlighting}[]
\SpecialCharTok{\textgreater{}} \FunctionTok{help}\NormalTok{(}\StringTok{"[["}\NormalTok{)}
\end{Highlighting}
\end{Shaded}

\subsubsection{4.3. Ayuda en formato HTML
navegable}\label{ayuda-en-formato-html-navegable}

En la mayoría de las instalaciones de R, la ayuda está disponible en
formato HTML mediante:

\begin{Shaded}
\begin{Highlighting}[]
\SpecialCharTok{\textgreater{}} \FunctionTok{help.start}\NormalTok{()}
\end{Highlighting}
\end{Shaded}

Esto lanza un navegador web que permite explorar las páginas de ayuda
mediante hipervínculos. El enlace ``Search Engine and Keywords'' en la
página cargada por \texttt{help.start()} es particularmente útil, ya que
contiene una lista de conceptos de alto nivel que permite buscar entre
las funciones disponibles --- una forma muy eficiente de orientarse
rápidamente y entender la amplitud de lo que R ofrece.

\subsubsection{4.4. Búsqueda de ayuda por palabra
clave}\label{buxfasqueda-de-ayuda-por-palabra-clave}

El comando \texttt{help.search} (alternativamente \texttt{??}) permite
buscar ayuda de varias formas:

\begin{Shaded}
\begin{Highlighting}[]
\SpecialCharTok{\textgreater{}}\NormalTok{ ??solve}
\end{Highlighting}
\end{Shaded}

Para más detalles y ejemplos, se recomienda consultar
\texttt{?help.search}.

\subsubsection{4.5. Ejecutar los ejemplos de un tema de
ayuda}\label{ejecutar-los-ejemplos-de-un-tema-de-ayuda}

Los ejemplos incluidos en un tema de ayuda normalmente pueden ejecutarse
con:

\begin{Shaded}
\begin{Highlighting}[]
\SpecialCharTok{\textgreater{}} \FunctionTok{example}\NormalTok{(topic)}
\end{Highlighting}
\end{Shaded}

\begin{center}\rule{0.5\linewidth}{0.5pt}\end{center}

\subsection{5. Sintaxis básica: comandos, sensibilidad a mayúsculas,
comentarios}\label{sintaxis-buxe1sica-comandos-sensibilidad-a-mayuxfasculas-comentarios}

Técnicamente, R es un \textbf{lenguaje de expresiones} (\emph{expression
language}) con una sintaxis muy simple.

\subsubsection{5.1. Sensibilidad a mayúsculas y
minúsculas}\label{sensibilidad-a-mayuxfasculas-y-minuxfasculas}

R \textbf{distingue mayúsculas de minúsculas} (\emph{case sensitive}),
como la mayoría de los paquetes basados en UNIX. Esto significa que
\texttt{A} y \texttt{a} son símbolos distintos y se refieren a variables
diferentes.

\subsubsection{5.2. Nombres válidos}\label{nombres-vuxe1lidos}

El conjunto de símbolos que pueden usarse en nombres de R depende del
sistema operativo y del país en el que se ejecuta R (técnicamente, de la
configuración regional o \emph{locale} en uso). Normalmente se permiten
todos los símbolos alfanuméricos (y en algunos países esto incluye
letras acentuadas), además del punto (\texttt{.}) y el guion bajo
(\texttt{\_}), con la restricción de que un nombre debe empezar con un
punto o una letra, y si empieza con un punto, el segundo carácter no
puede ser un dígito. Los nombres son, en la práctica, ilimitados en
longitud.

\begin{quote}
\textbf{Nota de portabilidad}: para código R portable (incluido el
destinado a paquetes de R), el manual recomienda usar únicamente los
caracteres \texttt{A–Za–z0–9}.
\end{quote}

\subsubsection{5.3. Expresiones y
asignaciones}\label{expresiones-y-asignaciones}

Los comandos elementales consisten en \textbf{expresiones} o
\textbf{asignaciones}:

\begin{itemize}
\tightlist
\item
  Si se da una expresión como comando, esta se \textbf{evalúa}, se
  \textbf{imprime} (a menos que se haga explícitamente invisible), y el
  valor se pierde.
\item
  Una asignación también evalúa una expresión y pasa el valor a una
  variable, pero el resultado \textbf{no se imprime automáticamente}.
\end{itemize}

\subsubsection{5.4. Separación de comandos y
agrupación}\label{separaciuxf3n-de-comandos-y-agrupaciuxf3n}

Los comandos se separan mediante punto y coma (\texttt{;}) o mediante un
salto de línea. Los comandos elementales pueden agruparse en una sola
expresión compuesta mediante llaves (\texttt{\{} y \texttt{\}}).

\subsubsection{5.5. Comentarios}\label{comentarios}

Los comentarios pueden colocarse casi en cualquier parte, comenzando con
el símbolo numeral (\texttt{\#}); todo lo que sigue hasta el final de la
línea es un comentario:

\begin{Shaded}
\begin{Highlighting}[]
\CommentTok{\# esto es un comentario completo}
\NormalTok{x }\OtherTok{\textless{}{-}} \DecValTok{5}  \CommentTok{\# esto es un comentario al final de una línea}
\end{Highlighting}
\end{Shaded}

\subsubsection{5.6. Límite de longitud de
línea}\label{luxedmite-de-longitud-de-luxednea}

Las líneas de comando introducidas en la consola están limitadas a
aproximadamente 4095 bytes (no caracteres).

\begin{center}\rule{0.5\linewidth}{0.5pt}\end{center}

\subsection{6. Asignación de
variables}\label{asignaciuxf3n-de-variables}

El operador de asignación estándar y tradicional en R es la flecha hacia
la izquierda:

\begin{Shaded}
\begin{Highlighting}[]
\SpecialCharTok{\textgreater{}}\NormalTok{ x }\OtherTok{\textless{}{-}} \FunctionTok{c}\NormalTok{(}\FloatTok{10.4}\NormalTok{, }\FloatTok{5.6}\NormalTok{, }\FloatTok{3.1}\NormalTok{, }\FloatTok{6.4}\NormalTok{, }\FloatTok{21.7}\NormalTok{)}
\end{Highlighting}
\end{Shaded}

Esto crea un vector llamado \texttt{x} con cinco elementos. El símbolo
\texttt{\textless{}-} puede leerse como ``se le asigna a'', y es la
forma documentada en el manual oficial para crear y nombrar objetos en
R. (El operador \texttt{=} también funciona como asignación en la
mayoría de los contextos, y es ampliamente usado en la práctica moderna,
pero el manual de referencia introduce y emplea consistentemente
\texttt{\textless{}-}).

\begin{center}\rule{0.5\linewidth}{0.5pt}\end{center}

\subsection{7. Historial de comandos}\label{historial-de-comandos}

Bajo muchas versiones de UNIX y en Windows, R ofrece un mecanismo para
recordar y reejecutar comandos previos. Las flechas verticales del
teclado permiten desplazarse hacia adelante y atrás por un
\textbf{historial de comandos} (\emph{command history}). Una vez
localizado un comando de esta forma, el cursor puede moverse dentro del
comando usando las flechas horizontales, y los caracteres pueden
eliminarse con la tecla \texttt{DEL} o añadirse con las demás teclas.

Las capacidades de recuperación y edición bajo UNIX son altamente
personalizables, mediante la biblioteca \texttt{readline}.
Alternativamente, el editor de texto Emacs ofrece mecanismos de soporte
más generales (vía ESS, \emph{Emacs Speaks Statistics}) para trabajar
interactivamente con R.

\begin{center}\rule{0.5\linewidth}{0.5pt}\end{center}

\subsection{8. Ejecutar comandos desde un archivo y redirigir
salida}\label{ejecutar-comandos-desde-un-archivo-y-redirigir-salida}

\subsubsection{8.1. Ejecutar un script}\label{ejecutar-un-script}

Si los comandos están almacenados en un archivo externo, por ejemplo
\texttt{commands.R} en el directorio de trabajo \texttt{work}, pueden
ejecutarse en cualquier momento de una sesión de R con el comando:

\begin{Shaded}
\begin{Highlighting}[]
\SpecialCharTok{\textgreater{}} \FunctionTok{source}\NormalTok{(}\StringTok{"commands.R"}\NormalTok{)}
\end{Highlighting}
\end{Shaded}

En Windows, la opción \textbf{Source} también está disponible en el menú
\textbf{File}.

\subsubsection{8.2. Redirigir la salida a un
archivo}\label{redirigir-la-salida-a-un-archivo}

La función \texttt{sink()} desvía toda la salida posterior de la consola
hacia un archivo externo:

\begin{Shaded}
\begin{Highlighting}[]
\SpecialCharTok{\textgreater{}} \FunctionTok{sink}\NormalTok{(}\StringTok{"record.lis"}\NormalTok{)}
\end{Highlighting}
\end{Shaded}

El comando, sin argumentos, restaura la salida a la consola nuevamente:

\begin{Shaded}
\begin{Highlighting}[]
\SpecialCharTok{\textgreater{}} \FunctionTok{sink}\NormalTok{()}
\end{Highlighting}
\end{Shaded}

\begin{center}\rule{0.5\linewidth}{0.5pt}\end{center}

\subsection{9. Permanencia de datos y el espacio de
trabajo}\label{permanencia-de-datos-y-el-espacio-de-trabajo}

Las entidades que R crea y manipula se conocen como \textbf{objetos}.
Pueden ser variables, arreglos de números, cadenas de caracteres,
funciones, o estructuras más generales construidas a partir de tales
componentes.

\subsubsection{9.1. Listar los objetos
actuales}\label{listar-los-objetos-actuales}

Durante una sesión de R, los objetos se crean y almacenan por nombre. El
comando:

\begin{Shaded}
\begin{Highlighting}[]
\SpecialCharTok{\textgreater{}} \FunctionTok{objects}\NormalTok{()}
\end{Highlighting}
\end{Shaded}

(alternativamente, \texttt{ls()}) muestra los nombres de (la mayoría de)
los objetos actualmente almacenados en R. La colección de objetos
actualmente almacenada se llama el \textbf{espacio de trabajo}
(\emph{workspace}).

\subsubsection{9.2. Eliminar objetos}\label{eliminar-objetos}

Para eliminar objetos está disponible la función \texttt{rm}:

\begin{Shaded}
\begin{Highlighting}[]
\SpecialCharTok{\textgreater{}} \FunctionTok{rm}\NormalTok{(x, y, z, ink, junk, temp, foo, bar)}
\end{Highlighting}
\end{Shaded}

\subsubsection{9.3. Guardado del espacio de trabajo entre
sesiones}\label{guardado-del-espacio-de-trabajo-entre-sesiones}

Todos los objetos creados durante una sesión de R pueden almacenarse
permanentemente en un archivo para uso en sesiones futuras. Al final de
cada sesión de R, se ofrece la oportunidad de guardar todos los objetos
actualmente disponibles. Si se indica que se desea hacerlo, los objetos
se escriben a un archivo llamado \texttt{.RData} en el directorio
actual, y las líneas de comando usadas en la sesión se guardan en un
archivo llamado \texttt{.Rhistory}.

Cuando R se inicia más tarde desde el mismo directorio, recarga el
espacio de trabajo desde ese archivo. Al mismo tiempo se recarga el
historial de comandos asociado.

\subsubsection{9.4. Recomendación oficial sobre directorios de
trabajo}\label{recomendaciuxf3n-oficial-sobre-directorios-de-trabajo}

El manual recomienda explícitamente usar \textbf{directorios de trabajo
separados} para distintos análisis realizados con R. Es bastante común
que se creen objetos con nombres como \texttt{x} e \texttt{y} durante un
análisis; estos nombres suelen ser significativos en el contexto de un
único análisis, pero puede ser bastante difícil decidir qué
representaban cuando se han realizado varios análisis distintos en el
mismo directorio.

\begin{quote}
\textbf{Nota práctica}: este es, precisamente, el problema que
\textbf{RStudio Projects} (y, para la gestión de paquetes por proyecto,
\texttt{renv}) resuelven de forma moderna y estructurada.
\end{quote}

\begin{center}\rule{0.5\linewidth}{0.5pt}\end{center}

\subsection{10. Vectores: la estructura
fundamental}\label{vectores-la-estructura-fundamental}

R opera sobre estructuras de datos con nombre. La estructura más simple
es el \textbf{vector numérico}, una entidad única que consiste en una
colección ordenada de números.

\subsubsection{\texorpdfstring{10.1. Crear un vector con
\texttt{c()}}{10.1. Crear un vector con c()}}\label{crear-un-vector-con-c}

Para configurar un vector llamado \texttt{x}, por ejemplo, consistente
en cinco números (10.4, 5.6, 3.1, 6.4 y 21.7), se usa la función
\texttt{c()} (\emph{combine}):

\begin{Shaded}
\begin{Highlighting}[]
\SpecialCharTok{\textgreater{}}\NormalTok{ x }\OtherTok{\textless{}{-}} \FunctionTok{c}\NormalTok{(}\FloatTok{10.4}\NormalTok{, }\FloatTok{5.6}\NormalTok{, }\FloatTok{3.1}\NormalTok{, }\FloatTok{6.4}\NormalTok{, }\FloatTok{21.7}\NormalTok{)}
\end{Highlighting}
\end{Shaded}

Esta es, según el manual, una \textbf{asignación de vector} usando la
función \texttt{c()}, que en este contexto puede tomar un número
arbitrario de argumentos vectoriales y cuyo valor resultante es un
vector obtenido al concatenar sus argumentos uno tras otro.

\subsubsection{10.2. Asignaciones
equivalentes}\label{asignaciones-equivalentes}

El manual señala explícitamente que las siguientes formas son
equivalentes a la anterior:

\begin{Shaded}
\begin{Highlighting}[]
\SpecialCharTok{\textgreater{}} \FunctionTok{assign}\NormalTok{(}\StringTok{"x"}\NormalTok{, }\FunctionTok{c}\NormalTok{(}\FloatTok{10.4}\NormalTok{, }\FloatTok{5.6}\NormalTok{, }\FloatTok{3.1}\NormalTok{, }\FloatTok{6.4}\NormalTok{, }\FloatTok{21.7}\NormalTok{))}
\end{Highlighting}
\end{Shaded}

y, en orden inverso (aunque poco usada en la práctica):

\begin{Shaded}
\begin{Highlighting}[]
\SpecialCharTok{\textgreater{}} \FunctionTok{c}\NormalTok{(}\FloatTok{10.4}\NormalTok{, }\FloatTok{5.6}\NormalTok{, }\FloatTok{3.1}\NormalTok{, }\FloatTok{6.4}\NormalTok{, }\FloatTok{21.7}\NormalTok{) }\OtherTok{{-}\textgreater{}}\NormalTok{ x}
\end{Highlighting}
\end{Shaded}

\begin{center}\rule{0.5\linewidth}{0.5pt}\end{center}

\subsection{11. Aritmética vectorizada y la regla de
reciclaje}\label{aritmuxe9tica-vectorizada-y-la-regla-de-reciclaje}

R opera sobre vectores completos sin necesidad de bucles explícitos para
operaciones aritméticas básicas. Por ejemplo, si \texttt{x} es un vector
ya definido, este comando crea un nuevo vector \texttt{y} con 11
entradas:

\begin{Shaded}
\begin{Highlighting}[]
\SpecialCharTok{\textgreater{}}\NormalTok{ y }\OtherTok{\textless{}{-}} \FunctionTok{c}\NormalTok{(x, }\DecValTok{0}\NormalTok{, x)}
\end{Highlighting}
\end{Shaded}

Los vectores pueden usarse en expresiones aritméticas, en cuyo caso las
operaciones se realizan \textbf{elemento por elemento}. Los vectores que
ocurren en la misma expresión no necesitan tener la misma longitud. Si
no tienen la misma longitud, el valor del resultado es el de la
expresión más larga, y los valores en el vector más corto se
\textbf{reciclan} tantas veces como sea necesario (posiblemente de forma
fraccionaria) hasta que coincidan en longitud con el vector más largo.
En particular, una constante numérica simple es, en este sentido, un
vector de longitud uno que se recicla tantas veces como sea necesario.
Esta es la conocida \textbf{regla de reciclaje} (\emph{recycling rule}).

Los operadores aritméticos elementales son los habituales \texttt{+},
\texttt{-}, \texttt{*}, \texttt{/} y \texttt{\^{}} (potencia). Además,
todas las funciones matemáticas comunes están disponibles: \texttt{log},
\texttt{exp}, \texttt{sin}, \texttt{cos}, \texttt{tan}, \texttt{sqrt},
entre otras.

\begin{center}\rule{0.5\linewidth}{0.5pt}\end{center}

\subsection{12. Generación de
secuencias}\label{generaciuxf3n-de-secuencias}

R provee varias formas de generar secuencias de números.

\subsubsection{\texorpdfstring{12.1. El operador
\texttt{:}}{12.1. El operador :}}\label{el-operador}

\begin{Shaded}
\begin{Highlighting}[]
\SpecialCharTok{\textgreater{}} \DecValTok{1}\SpecialCharTok{:}\DecValTok{30}
\end{Highlighting}
\end{Shaded}

genera la secuencia de enteros del 1 al 30.

\subsubsection{\texorpdfstring{12.2. La función
\texttt{seq()}}{12.2. La función seq()}}\label{la-funciuxf3n-seq}

La función \texttt{seq()} es una forma más general de generar
secuencias, permitiendo especificar incrementos distintos de 1, o la
longitud deseada de la secuencia en vez del valor final.

\subsubsection{\texorpdfstring{12.3. La función
\texttt{rep()}}{12.3. La función rep()}}\label{la-funciuxf3n-rep}

La función \texttt{rep()} permite generar secuencias repitiendo un
vector dado un número determinado de veces.

\begin{center}\rule{0.5\linewidth}{0.5pt}\end{center}

\subsection{13. Vectores lógicos}\label{vectores-luxf3gicos}

Además de las cantidades numéricas, R permite manipular
\textbf{cantidades lógicas} (\emph{logical quantities}). Los elementos
de un vector lógico pueden tomar los valores \texttt{TRUE},
\texttt{FALSE}, o \texttt{NA} (valor ``no disponible'').

Los vectores lógicos se generan mediante condiciones. Por ejemplo:

\begin{Shaded}
\begin{Highlighting}[]
\SpecialCharTok{\textgreater{}}\NormalTok{ temp }\OtherTok{\textless{}{-}}\NormalTok{ x }\SpecialCharTok{\textgreater{}} \DecValTok{13}
\end{Highlighting}
\end{Shaded}

establece \texttt{temp} como un vector de la misma longitud que
\texttt{x}, con valores \texttt{TRUE} o \texttt{FALSE} según si el
elemento correspondiente de \texttt{x} es mayor que 13 o no.

Los operadores condicionales habituales son \texttt{\textless{}},
\texttt{\textless{}=}, \texttt{\textgreater{}},
\texttt{\textgreater{}=}, \texttt{==} (igualdad exacta) y \texttt{!=}
(diferente). Además, si \texttt{c1} y \texttt{c2} son expresiones
lógicas, \texttt{c1\ \&\ c2} es su intersección (``y''),
\texttt{c1\ \textbar{}\ c2} es su unión (``o''), y \texttt{!c1} es la
negación de \texttt{c1}.

Los valores lógicos se coercionan a valores numéricos en expresiones
numéricas, con \texttt{FALSE} convertido en 0 y \texttt{TRUE} convertido
en 1.

\begin{center}\rule{0.5\linewidth}{0.5pt}\end{center}

\subsection{\texorpdfstring{14. Valores faltantes: \texttt{NA} y
\texttt{NaN}}{14. Valores faltantes: NA y NaN}}\label{valores-faltantes-na-y-nan}

En algunos casos, los componentes de un vector pueden no estar
disponibles. Cuando un elemento o valor no está disponible, se le asigna
un valor particular, \texttt{NA}. En general, cualquier operación sobre
un \texttt{NA} se convierte en \texttt{NA}. Esto permite que cualquier
proceso computacional que reciba \texttt{NA} como argumento se propague
de forma transparente a través de cálculos posteriores, dejando un
resultado igualmente desconocido o faltante.

La función \texttt{is.na(x)} da un resultado lógico que indica si cada
elemento de \texttt{x} es \texttt{NA} o no.

\subsubsection{\texorpdfstring{14.1. \texttt{NaN} --- Not a
Number}{14.1. NaN --- Not a Number}}\label{nan-not-a-number}

Hay una segunda clase de valores no disponibles producidos por cálculos
numéricos, conocida como valores \textbf{\texttt{NaN}} (\emph{Not a
Number}), por ejemplo el resultado de \texttt{0/0}. Los valores
\texttt{NaN} son todos \texttt{NA}, pero la afirmación inversa no es
cierta. La función \texttt{is.nan(x)} se usa específicamente para
detectar valores \texttt{NaN}.

\begin{center}\rule{0.5\linewidth}{0.5pt}\end{center}

\subsection{15. Vectores de caracteres}\label{vectores-de-caracteres}

Las secuencias de caracteres también son una estructura de vectores de
R, generalmente utilizadas como etiquetas para componentes de un objeto.
Por ejemplo:

\begin{Shaded}
\begin{Highlighting}[]
\SpecialCharTok{\textgreater{}}\NormalTok{ labs }\OtherTok{\textless{}{-}} \FunctionTok{paste}\NormalTok{(}\FunctionTok{c}\NormalTok{(}\StringTok{"X"}\NormalTok{,}\StringTok{"Y"}\NormalTok{), }\DecValTok{1}\SpecialCharTok{:}\DecValTok{10}\NormalTok{, }\AttributeTok{sep=}\StringTok{""}\NormalTok{)}
\end{Highlighting}
\end{Shaded}

crea un vector de longitud 10 mediante concatenación de pares de
cadenas. La función \texttt{paste()} permite combinar cualquier número
de argumentos, de manera elemento por elemento, en una sola cadena de
caracteres.

Las cadenas de caracteres se escriben usando comillas dobles
(preferencia indicada por el propio manual) o simples; cualquiera de las
dos puede usarse para ``escapar'' a la otra, como en
\texttt{"It\textquotesingle{}s\ important"}.

\begin{center}\rule{0.5\linewidth}{0.5pt}\end{center}

\subsection{16. Vectores de índices: seleccionar y modificar
subconjuntos}\label{vectores-de-uxedndices-seleccionar-y-modificar-subconjuntos}

Los subconjuntos de los elementos de un vector pueden seleccionarse
añadiendo un \textbf{vector de índices} entre corchetes
(\texttt{{[}{]}}), justo después del nombre del vector. El manual
oficial distingue cuatro tipos de vectores de índices:

\subsubsection{16.1. Vector lógico}\label{vector-luxf3gico}

El vector de índices es en sí mismo un vector lógico, de la misma
longitud que el vector del que se extraen los elementos. Los valores
correspondientes a \texttt{TRUE} en el vector de índices se seleccionan,
y los correspondientes a \texttt{FALSE} se omiten. Por ejemplo:

\begin{Shaded}
\begin{Highlighting}[]
\SpecialCharTok{\textgreater{}}\NormalTok{ y }\OtherTok{\textless{}{-}}\NormalTok{ x[}\SpecialCharTok{!}\FunctionTok{is.na}\NormalTok{(x)]}
\end{Highlighting}
\end{Shaded}

crea un objeto \texttt{y} que contiene los elementos no faltantes de
\texttt{x}, en el mismo orden en que ocurren en \texttt{x}.

\subsubsection{16.2. Vector de enteros
positivos}\label{vector-de-enteros-positivos}

El vector de índices puede consistir en un conjunto de valores enteros
positivos (o vacío). En este caso, los valores del vector de índices
deben estar en el conjunto \texttt{\{1,\ 2,\ ...,\ length(x)\}}. Los
elementos del vector seleccionados se especifican mediante su índice, en
el orden indicado.

\subsubsection{16.3. Vector de enteros
negativos}\label{vector-de-enteros-negativos}

El vector de índices puede consistir en un conjunto de números enteros
negativos. En este caso, el vector de índices especifica los valores que
\textbf{se excluyen}, en lugar de los que se incluyen.

\subsubsection{16.4. Vector de cadenas de
caracteres}\label{vector-de-cadenas-de-caracteres}

Finalmente, en el caso de un vector con nombres asociados a sus
elementos (mediante el atributo \texttt{names}), el vector de índices
puede ser un vector de cadenas de caracteres, refiriéndose a los
elementos por nombre en lugar de por posición.

\subsubsection{16.5. Reemplazo de elementos mediante
indexación}\label{reemplazo-de-elementos-mediante-indexaciuxf3n}

Un vector indexado puede usarse en el lado izquierdo de una asignación,
en cuyo caso la asignación solo opera sobre los elementos del vector
indicados por el vector de índices. Por ejemplo:

\begin{Shaded}
\begin{Highlighting}[]
\SpecialCharTok{\textgreater{}}\NormalTok{ x[}\FunctionTok{is.na}\NormalTok{(x)] }\OtherTok{\textless{}{-}} \DecValTok{0}
\end{Highlighting}
\end{Shaded}

reemplaza todos los valores \texttt{NA} de \texttt{x} por ceros.

\begin{center}\rule{0.5\linewidth}{0.5pt}\end{center}

\subsection{17. Objetos: modo, longitud y
atributos}\label{objetos-modo-longitud-y-atributos}

\subsubsection{\texorpdfstring{17.1. Modo
(\emph{mode})}{17.1. Modo (mode)}}\label{modo-mode}

Las entidades sobre las que opera R se conocen técnicamente como
\textbf{objetos}: ejemplos son vectores de valores numéricos (reales) o
complejos, vectores de valores lógicos y vectores de cadenas de
caracteres. Se conocen como estructuras ``atómicas'' porque todos sus
componentes son del mismo tipo, o \textbf{modo}, a saber:
\emph{numeric}, \emph{complex}, \emph{logical}, \emph{character} y
\emph{raw}.

Los vectores deben tener todos sus valores \textbf{del mismo modo}. Así,
cualquier vector dado debe ser inequívocamente lógico, numérico,
complejo, de caracteres, o raw. Una excepción aparente a esta regla es
el valor especial \texttt{NA}, para cantidades no disponibles (aunque,
de hecho, existen varios tipos de \texttt{NA}). Un vector puede estar
vacío y aun así tener un modo: el vector vacío de caracteres se lista
como \texttt{character(0)}, y el vector numérico vacío como
\texttt{numeric(0)}.

R también opera sobre objetos llamados \textbf{listas} (\emph{lists}),
de modo \emph{list}: son secuencias ordenadas de objetos que
individualmente pueden ser de cualquier modo. Las listas se conocen como
estructuras ``recursivas'' en lugar de ``atómicas'', porque sus
componentes pueden a su vez ser listas en sí mismas.

Las otras estructuras recursivas son las de modo \emph{function} y
\emph{expression}.

\subsubsection{\texorpdfstring{17.2. Longitud
(\emph{length})}{17.2. Longitud (length)}}\label{longitud-length}

Otra propiedad de todo objeto es su \textbf{longitud}. Las funciones
\texttt{mode(objeto)} y \texttt{length(objeto)} permiten averiguar el
modo y la longitud de cualquier estructura definida. Modo y longitud se
llaman, por esto, ``atributos intrínsecos'' de un objeto.

\subsubsection{17.3. Cambiar el tamaño de un
objeto}\label{cambiar-el-tamauxf1o-de-un-objeto}

Un objeto ``vacío'' puede aun así tener un modo. Por ejemplo:

\begin{Shaded}
\begin{Highlighting}[]
\SpecialCharTok{\textgreater{}}\NormalTok{ e }\OtherTok{\textless{}{-}} \FunctionTok{numeric}\NormalTok{()}
\end{Highlighting}
\end{Shaded}

hace de \texttt{e} un vector vacío de modo numérico. Una vez creado un
objeto de cualquier tamaño, pueden añadírsele nuevos componentes
simplemente asignándole un valor de índice fuera de su rango previo:

\begin{Shaded}
\begin{Highlighting}[]
\SpecialCharTok{\textgreater{}}\NormalTok{ e[}\DecValTok{3}\NormalTok{] }\OtherTok{\textless{}{-}} \DecValTok{17}
\end{Highlighting}
\end{Shaded}

esto convierte a \texttt{e} en un vector de longitud 3 (los dos primeros
componentes son, en este punto, ambos \texttt{NA}).

Inversamente, truncar el tamaño de un objeto solo requiere una
asignación apropiada:

\begin{Shaded}
\begin{Highlighting}[]
\SpecialCharTok{\textgreater{}}\NormalTok{ alpha }\OtherTok{\textless{}{-}}\NormalTok{ alpha[}\DecValTok{2} \SpecialCharTok{*} \DecValTok{1}\SpecialCharTok{:}\DecValTok{5}\NormalTok{]}
\end{Highlighting}
\end{Shaded}

o, más directamente:

\begin{Shaded}
\begin{Highlighting}[]
\SpecialCharTok{\textgreater{}} \FunctionTok{length}\NormalTok{(alpha) }\OtherTok{\textless{}{-}} \DecValTok{3}
\end{Highlighting}
\end{Shaded}

\subsubsection{17.4. Obtener y establecer
atributos}\label{obtener-y-establecer-atributos}

La función \texttt{attributes(objeto)} devuelve una lista de todos los
atributos no intrínsecos actualmente definidos para ese objeto. La
función \texttt{attr(objeto,\ nombre)} permite seleccionar un atributo
específico. Por ejemplo:

\begin{Shaded}
\begin{Highlighting}[]
\SpecialCharTok{\textgreater{}} \FunctionTok{attr}\NormalTok{(z, }\StringTok{"dim"}\NormalTok{) }\OtherTok{\textless{}{-}} \FunctionTok{c}\NormalTok{(}\DecValTok{10}\NormalTok{,}\DecValTok{10}\NormalTok{)}
\end{Highlighting}
\end{Shaded}

permite a R tratar a \texttt{z} como si fuera una matriz de 10 por 10.

\begin{center}\rule{0.5\linewidth}{0.5pt}\end{center}

\subsection{18. Coerción entre modos}\label{coerciuxf3n-entre-modos}

R permite cambios de modo prácticamente en cualquier lugar donde resulte
razonable hacerlo (y en algunos donde podría no serlo). Por ejemplo,
con:

\begin{Shaded}
\begin{Highlighting}[]
\SpecialCharTok{\textgreater{}}\NormalTok{ z }\OtherTok{\textless{}{-}} \DecValTok{0}\SpecialCharTok{:}\DecValTok{9}
\end{Highlighting}
\end{Shaded}

se podría escribir:

\begin{Shaded}
\begin{Highlighting}[]
\SpecialCharTok{\textgreater{}}\NormalTok{ digits }\OtherTok{\textless{}{-}} \FunctionTok{as.character}\NormalTok{(z)}
\end{Highlighting}
\end{Shaded}

tras lo cual \texttt{digits} es el vector de caracteres
\texttt{c("0",\ "1",\ "2",\ ...,\ "9")}. Una coerción adicional
reconstruye el vector numérico de nuevo:

\begin{Shaded}
\begin{Highlighting}[]
\SpecialCharTok{\textgreater{}}\NormalTok{ d }\OtherTok{\textless{}{-}} \FunctionTok{as.integer}\NormalTok{(digits)}
\end{Highlighting}
\end{Shaded}

Existe una amplia colección de funciones de la forma \texttt{as.algo()}
para coercionar de un modo a otro, o para dotar a un objeto de algún
otro atributo que aún no posea.

\begin{center}\rule{0.5\linewidth}{0.5pt}\end{center}

\subsection{19. La clase de un objeto}\label{la-clase-de-un-objeto}

Todos los objetos en R tienen una \textbf{clase}, reportada por la
función \texttt{class()}. Para vectores simples, esta es simplemente el
modo (por ejemplo \texttt{"numeric"}, \texttt{"logical"},
\texttt{"character"} o \texttt{"list"}), pero \texttt{"matrix"},
\texttt{"array"}, \texttt{"factor"} y \texttt{"data.frame"} son otros
valores posibles.

La clase es un atributo especial que permite un estilo de programación
orientado a objetos en R. Por ejemplo, si un objeto tiene clase
\texttt{"data.frame"}, se imprimirá de cierta forma, la función
\texttt{plot()} lo mostrará gráficamente de cierta manera, y otras
funciones llamadas ``genéricas'', como \texttt{summary()}, reaccionarán
a él como argumento de una forma sensible a su clase.

Para remover temporalmente los efectos de la clase, se usa la función
\texttt{unclass()}.

\begin{center}\rule{0.5\linewidth}{0.5pt}\end{center}

\subsection{20. Estructuras de control: condicionales y
bucles}\label{estructuras-de-control-condicionales-y-bucles}

R, como lenguaje, incluye las estructuras de control habituales de un
lenguaje de programación moderno.

\subsubsection{\texorpdfstring{20.1. Condicional
\texttt{if}}{20.1. Condicional if}}\label{condicional-if}

\begin{Shaded}
\begin{Highlighting}[]
\ControlFlowTok{if}\NormalTok{ (condicion) \{}
\NormalTok{  expresion\_verdadera}
\NormalTok{\} }\ControlFlowTok{else}\NormalTok{ \{}
\NormalTok{  expresion\_falsa}
\NormalTok{\}}
\end{Highlighting}
\end{Shaded}

\subsubsection{\texorpdfstring{20.2. Bucle
\texttt{for}}{20.2. Bucle for}}\label{bucle-for}

\begin{Shaded}
\begin{Highlighting}[]
\ControlFlowTok{for}\NormalTok{ (nombre }\ControlFlowTok{in}\NormalTok{ vector) \{}
\NormalTok{  expresion}
\NormalTok{\}}
\end{Highlighting}
\end{Shaded}

\subsubsection{\texorpdfstring{20.3. Bucles \texttt{while} y
\texttt{repeat}}{20.3. Bucles while y repeat}}\label{bucles-while-y-repeat}

R también incluye \texttt{while} (que repite mientras una condición sea
verdadera) y \texttt{repeat} (que repite indefinidamente hasta un
\texttt{break} explícito). Las palabras clave \texttt{break} y
\texttt{next} permiten, respectivamente, salir de un bucle o saltar a la
siguiente iteración.

\begin{center}\rule{0.5\linewidth}{0.5pt}\end{center}

\subsection{21. Funciones propias}\label{funciones-propias}

Una de las grandes fortalezas de R, según el propio manual, es la
facilidad con la que el usuario puede escribir sus propias funciones. La
sintaxis básica es:

\begin{Shaded}
\begin{Highlighting}[]
\NormalTok{nombre\_funcion }\OtherTok{\textless{}{-}} \ControlFlowTok{function}\NormalTok{(arg1, arg2, ...) \{}
  \CommentTok{\# cuerpo de la función}
\NormalTok{  resultado}
\NormalTok{\}}
\end{Highlighting}
\end{Shaded}

El valor devuelto por una función es, por defecto, el valor de la última
expresión evaluada en su cuerpo, aunque puede hacerse explícito con
\texttt{return()}.

Las funciones en R son objetos de primera clase: pueden almacenarse en
variables, pasarse como argumentos a otras funciones, y devolverse como
resultado de otras funciones, igual que cualquier otro objeto de R.

\begin{center}\rule{0.5\linewidth}{0.5pt}\end{center}

\subsection{22. Paquetes: el sistema de extensión de
R}\label{paquetes-el-sistema-de-extensiuxf3n-de-r}

Todas las funciones y conjuntos de datos de R se almacenan en
\textbf{paquetes} (\emph{packages}). Solo cuando un paquete se carga, su
contenido está disponible. Esto se hace tanto por eficiencia (la lista
completa de funciones ocuparía más memoria y tomaría más tiempo de
búsqueda que un subconjunto), como para ayudar a los desarrolladores de
paquetes, que quedan protegidos de colisiones de nombres con otro
código.

\subsubsection{22.1. Ver qué paquetes están
instalados}\label{ver-quuxe9-paquetes-estuxe1n-instalados}

\begin{Shaded}
\begin{Highlighting}[]
\SpecialCharTok{\textgreater{}} \FunctionTok{library}\NormalTok{()}
\end{Highlighting}
\end{Shaded}

(sin argumentos) muestra los paquetes instalados en el sistema.

\subsubsection{22.2. Cargar un paquete}\label{cargar-un-paquete}

Para cargar un paquete particular, por ejemplo el paquete \textbf{boot}:

\begin{Shaded}
\begin{Highlighting}[]
\SpecialCharTok{\textgreater{}} \FunctionTok{library}\NormalTok{(boot)}
\end{Highlighting}
\end{Shaded}

\subsubsection{22.3. Instalar y actualizar
paquetes}\label{instalar-y-actualizar-paquetes}

Los usuarios conectados a Internet pueden usar las funciones
\texttt{install.packages()} y \texttt{update.packages()} (disponibles
también a través del menú \texttt{Packages} en las interfaces gráficas
de Windows y macOS) para instalar y actualizar paquetes.

\subsubsection{22.4. Ver qué paquetes están
cargados}\label{ver-quuxe9-paquetes-estuxe1n-cargados}

\begin{Shaded}
\begin{Highlighting}[]
\SpecialCharTok{\textgreater{}} \FunctionTok{search}\NormalTok{()}
\end{Highlighting}
\end{Shaded}

muestra la lista de búsqueda (\emph{search list}). Algunos paquetes
pueden estar cargados pero no aparecer en la lista de búsqueda (por usar
espacios de nombres); estos se incluyen en la lista entregada por:

\begin{Shaded}
\begin{Highlighting}[]
\SpecialCharTok{\textgreater{}} \FunctionTok{loadedNamespaces}\NormalTok{()}
\end{Highlighting}
\end{Shaded}

\subsubsection{22.5. Paquetes estándar (base) vs.~paquetes
contribuidos}\label{paquetes-estuxe1ndar-base-vs.-paquetes-contribuidos}

Los paquetes \textbf{estándar} (o \emph{base}) se consideran parte del
código fuente de R: contienen las funciones básicas que permiten que R
funcione, junto con los conjuntos de datos y las funciones estadísticas
y gráficas estándar descritas en el manual oficial. Deben estar
disponibles automáticamente en cualquier instalación de R.

Existen, además, miles de \textbf{paquetes contribuidos}, escritos por
muchos autores distintos. Algunos implementan métodos estadísticos
especializados, otros dan acceso a datos o hardware, y otros están
diseñados para complementar libros de texto. Algunos (los paquetes
``recomendados'') se distribuyen junto con toda distribución binaria de
R. La mayoría están disponibles para descarga desde CRAN
(cran.r-project.org y sus espejos) y otros repositorios como
Bioconductor (bioconductor.org).

\begin{center}\rule{0.5\linewidth}{0.5pt}\end{center}

\subsection{23. Espacios de nombres
(namespaces)}\label{espacios-de-nombres-namespaces}

Los paquetes tienen \textbf{espacios de nombres} (\emph{namespaces}),
que cumplen tres funciones:

\begin{enumerate}
\def\labelenumi{\arabic{enumi}.}
\tightlist
\item
  Permiten al autor de un paquete ocultar funciones y datos destinados
  únicamente a uso interno.
\item
  Evitan que las funciones se rompan cuando un usuario (u otro autor de
  paquete) elige un nombre que colisiona con uno ya existente en el
  paquete.
\item
  Proveen una forma de referirse a un objeto dentro de un paquete
  particular.
\end{enumerate}

\subsubsection{\texorpdfstring{23.1. El operador de doble dos puntos
\texttt{::}}{23.1. El operador de doble dos puntos ::}}\label{el-operador-de-doble-dos-puntos}

Por ejemplo, \texttt{t()} es la función de transposición en R, pero un
usuario podría definir su propia función llamada \texttt{t}. Los
espacios de nombres evitan que la definición del usuario tenga
precedencia y rompa cualquier función que intente transponer una matriz.
El operador \texttt{::} selecciona definiciones de un espacio de nombres
particular: la función de transposición siempre estará disponible como
\texttt{base::t}, porque está definida en el paquete \texttt{base}. Solo
las funciones exportadas desde un paquete pueden recuperarse de esta
forma.

\subsubsection{\texorpdfstring{23.2. El operador de triple dos puntos
\texttt{:::}}{23.2. El operador de triple dos puntos :::}}\label{el-operador-de-triple-dos-puntos}

El operador \texttt{:::} actúa como el operador \texttt{::}, pero
también permite acceder a objetos ocultos. Es más probable que los
usuarios recurran a la función \texttt{getAnywhere()}, que busca en
múltiples paquetes.

\begin{center}\rule{0.5\linewidth}{0.5pt}\end{center}

\subsection{24. Buenas prácticas según el manual
oficial}\label{buenas-pruxe1cticas-seguxfan-el-manual-oficial}

A modo de síntesis, estas son recomendaciones que el propio manual de
CRAN hace explícitas a lo largo del texto, antes de entrar en
manipulación de datos:

\begin{itemize}
\tightlist
\item
  Usar \textbf{un directorio de trabajo separado por análisis}, para
  evitar confusión entre objetos con el mismo nombre (típicamente
  \texttt{x}, \texttt{y}) creados en distintos análisis no relacionados.
\item
  Familiarizarse primero con \texttt{help()}, \texttt{?},
  \texttt{help.search()} y \texttt{help.start()} antes de buscar ayuda
  externa: el sistema de ayuda integrado de R es, según el manual,
  completo y debe ser el primer recurso.
\item
  Recordar que R \textbf{distingue mayúsculas de minúsculas} en todos
  los nombres.
\item
  Aprovechar la naturaleza \textbf{vectorizada} de R: operar sobre
  vectores completos en lugar de escribir bucles explícitos siempre que
  sea razonable, ya que esto es parte del diseño fundamental del
  lenguaje, no solo una optimización de estilo.
\item
  Entender la diferencia entre \texttt{NA} (valor no disponible, de
  cualquier tipo) y \texttt{NaN} (resultado numérico inválido
  específicamente, como \texttt{0/0}), y usar \texttt{is.na()} /
  \texttt{is.nan()} según corresponda.
\item
  Recordar que, en R, los resultados de un análisis estadístico
  normalmente se almacenan en un objeto para su examen posterior, en
  lugar de imprimirse exhaustivamente de inmediato.
\end{itemize}

\begin{center}\rule{0.5\linewidth}{0.5pt}\end{center}

\subsection{25. Referencias oficiales}\label{referencias-oficiales}

\begin{itemize}
\tightlist
\item
  \textbf{CRAN --- ``An Introduction to R'' (HTML, versión release)}:
  https://cran.r-project.org/doc/manuals/r-release/R-intro.html
\item
  \textbf{CRAN --- ``An Introduction to R'' (PDF, versión release)}:
  https://cran.r-project.org/doc/manuals/r-release/R-intro.pdf
\item
  \textbf{R Manuals (réplica oficial restilizada) --- Índice general}:
  https://rstudio.github.io/r-manuals/r-intro/index.html
\item
  \textbf{R Manuals --- Cap. 1: Introduction and preliminaries}:
  https://rstudio.github.io/r-manuals/r-intro/Introduction-and-preliminaries.html
\item
  \textbf{R Manuals --- Cap. 2: Simple manipulations; numbers and
  vectors}:
  https://rstudio.github.io/r-manuals/r-intro/Simple-manipulations-numbers-and-vectors.html
\item
  \textbf{R Manuals --- Cap. 3: Objects, their modes and attributes}:
  https://rstudio.github.io/r-manuals/r-intro/Objects.html
\item
  \textbf{R Manuals --- Cap. 9: Grouping, loops and conditional
  execution}:
  https://rstudio.github.io/r-manuals/r-intro/Loops-and-conditional-execution.html
\item
  \textbf{R Manuals --- Cap. 10: Writing your own functions}:
  https://rstudio.github.io/r-manuals/r-intro/Writing-your-own-functions.html
\item
  \textbf{R Manuals --- Cap. 13: Packages}:
  https://rstudio.github.io/r-manuals/r-intro/Packages.html
\item
  \textbf{CRAN --- R FAQ (oficial)}:
  https://cran.r-project.org/doc/manuals/r-release/R-FAQ.html
\item
  \textbf{CRAN --- Manuales completos de R (índice)}:
  https://cran.r-project.org/manuals.html
\end{itemize}

\begin{center}\rule{0.5\linewidth}{0.5pt}\end{center}

\textbf{Edison Achalma} Economista --- Universidad Nacional de San
Cristóbal de Huamanga (UNSCH) Ayacucho, Perú ORCID:
\href{https://orcid.org/0000-0001-6996-3364}{0000-0001-6996-3364}

\section{Publicaciones Similares}\label{publicaciones-similares}

Si te interesó este artículo, te recomendamos que explores otros blogs y
recursos relacionados que pueden ampliar tus conocimientos. Aquí te dejo
algunas sugerencias:

\begin{enumerate}
\def\labelenumi{\arabic{enumi}.}
\tightlist
\item
  \href{https://numerus-scriptum.netlify.app/r/2020-06-10-011-instalacion-de-r/index.pdf}{\faIcon{file-pdf}}
  \href{https://numerus-scriptum.netlify.app/r/2020-06-10-011-instalacion-de-r}{011
  Instalacion De R}
\item
  \href{https://numerus-scriptum.netlify.app/r/2020-06-10-012-configurar-entorno-virtual-r/index.pdf}{\faIcon{file-pdf}}
  \href{https://numerus-scriptum.netlify.app/r/2020-06-10-012-configurar-entorno-virtual-r}{012
  Configurar Entorno Virtual R}
\item
  \href{https://numerus-scriptum.netlify.app/r/2020-06-10-013-lo-que-debemos-saber-de-r/index.pdf}{\faIcon{file-pdf}}
  \href{https://numerus-scriptum.netlify.app/r/2020-06-10-013-lo-que-debemos-saber-de-r}{013
  Lo Que Debemos Saber De R}
\item
  \href{https://numerus-scriptum.netlify.app/r/2021-04-05-02-manipulacion-de-datos/index.pdf}{\faIcon{file-pdf}}
  \href{https://numerus-scriptum.netlify.app/r/2021-04-05-02-manipulacion-de-datos}{02
  Manipulacion De Datos}
\item
  \href{https://numerus-scriptum.netlify.app/r/2021-04-12-03-visualizacion-de-datos/index.pdf}{\faIcon{file-pdf}}
  \href{https://numerus-scriptum.netlify.app/r/2021-04-12-03-visualizacion-de-datos}{03
  Visualizacion De Datos}
\item
  \href{https://numerus-scriptum.netlify.app/r/2022-11-07-04-modelo-de-machine-learning-i-analisis-exploratorio/index.pdf}{\faIcon{file-pdf}}
  \href{https://numerus-scriptum.netlify.app/r/2022-11-07-04-modelo-de-machine-learning-i-analisis-exploratorio}{04
  Modelo De Machine Learning I Analisis Exploratorio}
\item
  \href{https://numerus-scriptum.netlify.app/r/2022-11-14-05-modelo-de-machine-learning-ii-modelo-de-clasificacion/index.pdf}{\faIcon{file-pdf}}
  \href{https://numerus-scriptum.netlify.app/r/2022-11-14-05-modelo-de-machine-learning-ii-modelo-de-clasificacion}{05
  Modelo De Machine Learning Ii Modelo De Clasificacion}
\item
  \href{https://numerus-scriptum.netlify.app/r/2022-11-21-06-modelo-de-machine-learning-iii-modelo-de-regresion/index.pdf}{\faIcon{file-pdf}}
  \href{https://numerus-scriptum.netlify.app/r/2022-11-21-06-modelo-de-machine-learning-iii-modelo-de-regresion}{06
  Modelo De Machine Learning Iii Modelo De Regresion}
\item
  \href{https://numerus-scriptum.netlify.app/r/2022-11-28-07-modelo-de-machine-learning-iv-tex-mining/index.pdf}{\faIcon{file-pdf}}
  \href{https://numerus-scriptum.netlify.app/r/2022-11-28-07-modelo-de-machine-learning-iv-tex-mining}{07
  Modelo De Machine Learning Iv Tex Mining}
\end{enumerate}

Esperamos que encuentres estas publicaciones igualmente interesantes y
útiles. ¡Disfruta de la lectura!






\end{document}
